% Revised manuscript, Section 2 -- draft for review (2026-10-24)
\section{Multiplicative Nulling in the Sedenion Algebra}

\subsection{Array model and algebra}
An 8-element CRPA sampled at baseband produces, at chip index $t$, a snapshot
$z(t)\in\mathbb{C}^8$. Writing $I_m,Q_m$ for the in-phase and quadrature outputs of element
$m$, we identify the snapshot with a sedenion
\begin{equation}
X(t)=\sum_{m=0}^{7}\bigl(I_m(t)\,e_m+Q_m(t)\,e_{m+8}\bigr)=\bigl(a(t),b(t)\bigr)\in
\mathbb{O}\oplus\mathbb{O}e_8=\mathbb{S},
\label{eq:embed}
\end{equation}
an isometry $\mathbb{C}^8\to\mathbb{R}^{16}\cong\mathbb{S}$. The received snapshot is
$X(t)=\sqrt{P_s}\,c(t-\tau_0)\,u_{\mathrm{sat}}+\sum_r\sqrt{P_r}\,j_r(t)\,u_r+N(t)$, with
$u_{\mathrm{sat}},u_r$ unit spatial signatures, $c(\cdot)$ the C/A code and $N$ white noise.
Sedenion multiplication is the Cayley--Dickson doubling
\begin{equation}
(a,b)\cdot(c,d)=\bigl(ac-d^{*}b,\;da+bc^{*}\bigr),
\label{eq:cd}
\end{equation}
bilinear over $\mathbb{R}$, so for fixed $v$ the map $X\mapsto X\cdot v$ is a real-linear
filter $L_v$ on $\mathbb{R}^{16}$. The embedding~\eqref{eq:embed} is a coordinate choice: it
does not by itself give the product a physical meaning, and the results below hold for any
fixed isometric embedding composed with the gauge of \S\ref{sec:gauge}.

\subsection{Zero divisors and the closed-form annihilator fibre}
For $u=(a,b)$ with $a,b\in\mathrm{Im}(\mathbb{O})$, $\|a\|=\|b\|=1$, $\langle a,b\rangle=0$
(the Moreno cone, diffeomorphic to the Stiefel manifold $V_2(\mathbb{R}^7)$
\cite{moreno1998}), let $W_{ab}=\mathrm{span}\{a,b,ab\}\subset\mathrm{Im}(\mathbb{O})$.

\begin{theorem}[Annihilator fibre \cite{agourakis2026pg32}]\label{thm:fibre}
For every $c\in W_{ab}^\perp$ (a 4-dimensional subspace),
\begin{equation}
v(c)=\tfrac{1}{\sqrt2}\bigl(c,\,(ac)b\bigr)\quad\text{satisfies}\quad u\cdot v(c)=0,
\label{eq:fibre}
\end{equation}
the map $c\mapsto v(c)$ is linear and norm-preserving ($\|v(c)\|=\|c\|$), and
$\dim\ker L(u)=4$.
\end{theorem}
The identity $u\cdot v(c)=0$ is formalised in Lean~4/Mathlib as \texttt{nullspace} and the
dimension count as \texttt{finrank\_ker\_Lu}. Computing $v(c)$ costs two octonion products.

\subsection{Anti-partner gauge}\label{sec:gauge}
A physical interferer signature does not lie on the Moreno cone, and the satellite signature
bears no particular relation to the fibre. Both are fixed by an orthogonal change of
coordinates. Let $u_{\mathrm{zd}}=\tfrac{1}{\sqrt2}(e_3+e_{12})$ be a canonical cone point and
$u_{\mathrm{anti}}=\tfrac{1}{\sqrt2}(e_3-e_{12})$ its anti-partner. Given an estimate
$\hat u_1$ of the interferer signature,
\begin{align}
T_1&=H(\hat u_1\mapsto u_{\mathrm{zd}}),\\
T_2&=H\bigl(\mathrm{unit}(T_1u_{\mathrm{sat}}^{\perp})\mapsto u_{\mathrm{anti}}\bigr),\qquad
u_{\mathrm{sat}}^{\perp}=u_{\mathrm{sat}}-\langle\hat u_1,u_{\mathrm{sat}}\rangle\hat u_1,
\end{align}
where $H(x\mapsto y)$ is the Householder reflection exchanging unit vectors $x,y$, and
$T=T_2T_1$. Because $T_1u_{\mathrm{sat}}^\perp\perp u_{\mathrm{zd}}$ and
$u_{\mathrm{anti}}\perp u_{\mathrm{zd}}$, the reflection $T_2$ fixes $u_{\mathrm{zd}}$, so
$T\hat u_1=u_{\mathrm{zd}}$ exactly while the satellite's interferer-orthogonal component is
carried onto $u_{\mathrm{anti}}$. The filter applied to the array is
\begin{equation}
y(X)=(TX)\cdot v(c),\qquad c\in W_{\mathrm{zd}}^\perp ,
\label{eq:filter}
\end{equation}
which annihilates $\hat u_1$ for every $c$ and passes the satellite with amplitude equal to
$\|u_{\mathrm{sat}}^\perp\|=\sqrt{1-\langle\hat u_1,u_{\mathrm{sat}}\rangle^2}$, the same
passband limit as any linear beamformer constrained to null $\hat u_1$ (\S5: $0.9895$
measured vs.\ $0.9899$ theoretical).

We stress what the gauge is: an orthogonal $16\times16$ map built from two outer products,
of the same form and cost as a rank-1 projector. It is the step that makes the algebraic
construction usable on physical signatures; it is not a source of advantage over linear
beamforming, and we do not claim one.

\subsection{Constraints inside the fibre}\label{sec:fibreconstraints}
Let $\{c_1,\dots,c_4\}$ be an orthonormal basis of $W_{\mathrm{zd}}^\perp$ and write
$c=\sum_k\alpha_kc_k$. Every readout of~\eqref{eq:filter} is linear in
$\alpha\in\mathbb{R}^4$. Choosing the quadrature pair $(e_i,e_j)$ of the downstream
correlator as I/Q, we impose:
\begin{itemize}
\item \emph{Phase lock}: $\langle e_j,\,y(u_{\mathrm{sat}})\rangle=0$, so the satellite enters
the correlator with zero carrier phase regardless of how $\hat u_1$ was re-estimated;
\item \emph{Second interferer}: $\langle e_i,y(\hat u_2)\rangle=\langle e_j,y(\hat u_2)\rangle=0$.
\end{itemize}
Stacking $m\le3$ such rows gives $L\alpha=0$ with $L\in\mathbb{R}^{m\times4}$; we take the unit
$\alpha$ in $\ker L$ maximising the in-phase satellite amplitude
$\langle e_i,y(u_{\mathrm{sat}})\rangle$. The solve is a $3\times4$ null-space computation,
comparable in cost to the $3\times3$ Gram solve of an LCMV beamformer with the same
constraints. The fibre offers four real degrees of freedom against roughly fifteen for an
LCMV weight in $\mathbb{C}^8$, which is consistent with the 4~dB dual-nulling gap reported
in \S5.

\subsection{Two-stage cascade}
To null a second interferer (or the tangent direction of a drifting one) across all output
components rather than only the I/Q readout, a second stage is applied to the residual:
\begin{equation}
y^{(2)}(X)=\bigl(T^{(2)}y^{(1)}(X)\bigr)\cdot v_2,\qquad
T^{(2)}\,\mathrm{unit}\bigl(y^{(1)}(\hat u_2)\bigr)=u_{\mathrm{zd}},
\label{eq:cascade}
\end{equation}
so that $y^{(2)}(\hat u_1)=y^{(2)}(\hat u_2)=0$ exactly. With $\hat u_1,\hat u_2$ the two
leading eigenvectors of the same sample covariance used by the baselines, the cascade is the
multiplicative counterpart of rank-2 orthogonal projection and performs at its level (\S5).

\subsection{Leakage bounds}\label{sec:bound}
\begin{theorem}[Single-stage leakage]\label{thm:leak}
If $\hat u\cdot v=0$, then for any perturbation $\delta\in\mathbb{S}$,
\begin{equation}
\|(\hat u+\delta)\cdot v\|^2\le 2\,\|\delta\|^2\|v\|^2 .
\end{equation}
\end{theorem}
\begin{proof}[Proof sketch]
By bilinearity, $(\hat u+\delta)\cdot v=\delta\cdot v$. Writing $\delta=(p,q)$, $v=(c,d)$ and
applying~\eqref{eq:cd}, each octonion component is a difference or sum of two octonion
products; the parallelogram inequality and the octonion composition law
$\|xy\|=\|x\|\|y\|$ give $\|\delta\cdot v\|^2\le 2(\|p\|^2+\|q\|^2)(\|c\|^2+\|d\|^2)$.
\end{proof}

\begin{theorem}[Cascade leakage]\label{thm:cascade}
If $(\hat u\cdot v_1)\cdot v_2=0$, then
$\|((\hat u+\delta)\cdot v_1)\cdot v_2\|^2\le 4\,\|\delta\|^2\|v_1\|^2\|v_2\|^2$.
\end{theorem}

Both theorems are formalised in Lean~4 (v4.15.0, Mathlib) as \texttt{leakage\_bound} and
\texttt{cascade\_leakage\_bound}, with no \texttt{sorry} and only the standard axioms
\texttt{propext}, \texttt{Classical.choice}, \texttt{Quot.sound}. Since $T$ is orthogonal, the
bounds apply unchanged to the gauged filter~\eqref{eq:filter}, with $\|\delta\|$ the
signature-estimation error.

For unit $v$, Theorem~\ref{thm:leak} gives a guaranteed null depth of at least
$-20\log_{10}\|\delta\|-3$~dB. This is the operational content of the result: null depth is
set by how well the interferer signature is known, through calibration and estimation, and not
by the algebra. It is the same dependence a linear signature-driven nuller has; what the
theorem adds is that for the multiplicative construction it is proven rather than assumed.
Exact signature knowledge ($\delta=0$) yields annihilation to the floating-point floor
(139~dB in FP32, 313~dB in FP64), which is the limit case of the bound and not a performance
figure. In \S5 the measured single-stage residual lies within a factor $1.4$ of the bound at
every tested $\|\delta\|$, with no violations in 1{,}792 trials.
